\chapter{Objectives}
\label{ch:objectives}

The aim of the project is to design and develop a gamified web-security training platform where students practise penetration testing on realistic vulnerable applications, with interactive learning, performance evaluation and instructor analytics.

The aim breaks into six objectives. Each closes one gap from Table~\ref{tab:gaps}, and Table~\ref{tab:objectives} gives the success criterion for each. The adaptive behaviour comes from three objectives working together: scoring produces a performance record (O3), a simple model turns that record into a prediction and a recommendation (O4), and the instructor sees the result (O5).

\begin{enumerate}[label=\textbf{O\arabic*},leftmargin=2.6em,itemsep=0.45em]
\item Develop an intentionally vulnerable web application with OWASP Top~10 vulnerabilities~\cite{OwaspTop10}, delivered as challenges inside the platform.
\item Design a gamification module with points, levels, badges and achievements, where each element is designed around the documented negative effects of gamification.
\item Develop an automated scoring mechanism that evaluates each attempt from completion, accuracy, time and hint use, and records every attempt as the learner's performance record.
\item Provide hints and learning resources during exercises, and use a simple, interpretable model over the performance record to estimate a learner's likely difficulty and recommend a suitably difficult next challenge.
\item Develop an instructor dashboard that shows progress, challenge completion and performance for each learner and for the cohort, with at-risk learners flagged early.
\item Evaluate the effect of the platform on learner engagement and practical cybersecurity skill.
\end{enumerate}

\begin{table}[!htbp]
\centering
\caption{Objectives, the gap each one closes, and its success criterion.}
\label{tab:objectives}
\tablefont
\begin{tabularx}{\linewidth}{@{}c P{3.3cm} c L@{}}
\toprule
\textbf{ID} & \textbf{Objective} & \textbf{Gap} & \textbf{Success criterion} \\
\midrule
O1 & Vulnerable web application & G1 & Challenges from several OWASP Top~10 families, each exploitable end to end in a contained lab, with a separate flag for each learner checked on the server. \\
O2 & Gamification module & G2 & Points, levels, badges and achievements, each with a design choice justified against the negative-effects literature, and no leaderboard. \\
O3 & Automated scoring & G3 & Every attempt scored automatically as a function of completion, accuracy and time, with no manual grading. \\
O4 & Adaptive support & G4 & A hint ladder on every challenge, and a next-challenge recommendation validated against held-out learner data. \\
O5 & Instructor dashboard & G5 & Per-learner and cohort progress, completion and scores from live data, with at-risk learners flagged. \\
O6 & Effectiveness evaluation & G6 & A pilot with a student cohort that reports the change in practical skill between a pre-test and a post-test, and a measured engagement outcome. \\
\bottomrule
\end{tabularx}
\end{table}

The objectives depend on one another, as Figure~\ref{fig:objectives} shows. Scoring needs gradable attempts from the lab. Gamification, adaptation and the dashboard all read the performance record that scoring writes, and the pilot can only run once the learner-facing parts are live. This order decided what was built first: O1 and O3 are complete, O2, O4 and O5 have working parts, and O6 is planned. Chapter~\ref{ch:results} gives the status of each objective in detail.

\diagramfig{objectives}{Dependencies between the objectives}{Status at mid-semester: a solid border means built, a dashed orange border means partly built, and a dashed grey border means planned.}{fig:objectives}
